\chapter{20}

Eine Gruppe Kinder spritzte mit Wasserpistolen die Hitze weg. Sie quietschten wie kleine Ferkel. Selbst einzelne Erwachsene begannen, sich verhalten zu bewegen. 
Schliesslich johlte einer beim Refrain lauthals mit. Er wirkte etwas neben der Spur. Einige
grinsten ihn an. Über ihnen zog ein Mauersegler seine Kreise. Louis hängte seinen
Blick der Flugbahn nach. Es duftete nach Sonnencrème. Unverwechselbar. Musik war wie Nachhausekommen. 
Wenigstens darauf konnte er zählen. 

Auf die Steigerung durch die Zweistimmigkeit war Louis nicht
gefasst. Die voluminöse Stimme der Frau, ergänzt durch den kraftvollen
Bass des Mannes, liess nun auch ihn in sich tanzen. Louis staunte. Wie
raffiniert sie den Song arrangiert hatten, kunstvoll adaptiert an ihre Instrumente
und ihre umwerfenden Stimmen. Wie sie sich dabei anschauten.

\sectionbreak 

\emph{«Ich werde auf dich warten. Und sollte ich zurückbleiben, warte auf 
mich. Wir schworen, wir reisen zusammen.»} 

\sectionbreak  

Die Worte stimmten Louis wehmütig.

Der Beifall nahm den Platz breit ein. Das Paar
verneigte sich tief und mehrmals. Die vielen Zweifrankenstücke,
Fünfliber und einzelnen Zehnernoten nickten sie dankend ab.

«Bis bald, unsere halbe Stunde ist um, wir müssen weiterziehen und
vielen Dank.»

Die Frau strahlte einhändig winkend in die Menge. Die beiden begannen,
ihre Instrumente einzupacken. Nach und nach schlenderten die Menschen
weiter. 
Louis ging näher auf sie zu. Er blieb verzaubert vor ihnen stehen.
Der junge Mann rollte in der Hocke vor Louis’ Füssen die Kabel ein.

«Das war gut,» sagte Louis begeistert, «hat mir sehr gefallen.»

Der Mann schaute hoch und strahlte. Sein Dialekt verriet den
Zürcher.

«Oh, danke, Springsteen ist ein lässiger Typ, kommt eigentlich immer gut
an.»

Während er den Gitarrenkoffer öffnete, beugte sich Louis zu ihm hinunter.

«Darf ich dich was fragen?»

«Ja, klar.»

Der Mann nickte, stand auf und reichte Louis seine heisse Hand. Er
schwitzte wie ein Bär.

«Bin Meo, hi, und du?»

«Ciro, freut mich, Meo.»

Louis schauderte. Zum ersten Mal stellte er sich mit neuem Namen vor.
Ciro Lorente. Ein gelungener Anfang. Mit ihrem Geigenkoffer in der Hand und einem schwarzen, grossen Rucksack am Rücken kam nun auch die Frau näher. Trotz des Gepäcks hatte sie einen federleichten Gang. Meo deutete mit einer Handbewegung zu Louis und blickte zwischen den beiden hin und her.

«Das ist Ciro,» und zu Louis gewandt, «Nessi, Vanessa, meine Freundin.»

Sie begrüssten sich. Die beiden schauten Louis aufmerksam an.
Er erzählte ihnen von einer gestohlenen Gitarre und einem weggekommenen
Handy. Von Wertsachen, die er dringend ersetzen müsse, von einem halben
Jahr Herumreisen und wie er von Mailand aus losgezogen sei. Er brauche
Zeit für sich. Wieder einmal in seiner alten Schweizerheimat. Er wolle
sich vorübergehend mit Strassenmusik durchbringen.
Dabei verwarf er die Hände wie ein Laienschauspieler.

«Und nun ohne Gitarre, ziemlich schlecht.»

Er seufzte.

«Hm.»

Meo wirkte betroffen. Vanessa verengte die Augen und streckte den
Kopf vor.

«Nun brauchst du eine Neue!?»

«Ja, dringend, das wird mit meiner Versicherung dauern,» Louis kam in Fahrt, 
«aber ich kann nicht warten, und als ich euch spielen hörte,
dachte ich, ich würde euch fragen, ob ihr mir einen Tipp habt. Ich sehe,
Meo,» und Louis zeigte auf Meos Gitarrenkoffer, «du hast ja ein
Prachtstück, wow, eine «Gibson», ja?»

«Genau.»

Es klang gedehnt, sichtlich stolz und Meo schloss den Koffer.

«Lohnt sich, mein Liebling hilft mit, unser Leben zu finanzieren.»

Louis stutzte.

«Die Gitarre, mein ich. Ich habe natürlich zwei Lieblinge und kann mich
manchmal schwer zwischen den beiden entscheiden - »

Meo lachte mit Seitenblick zu Vanessa. Sie boxte ihn sanft in die Seite und machte ein beleidigtes Gesicht. 
Louis schmunzelte.

«Du suchst also eine neue Gitarre, bist du Musiker?»

«Ja, und dabei müsste es schnell gehen. Sonst endet meine
Schweizerreise, bevor sie richtig angefangen hat.»

Jetzt lächelte Louis gequält.
Die beiden sahen ihn mitfühlend an. Meo schien nachzudenken.

«Ich kenne mich etwas aus mit Gitarren in Zürich,» Louis hängte sich ihm
gleich an die Lippen, «ich bin von Urdorf. Kommt natürlich darauf an,
welche Liga es denn sein soll und was du zahlen willst.»

«Ich suche wieder eine «Taylor», hatte die «Taylor Koa», die war
grossartig,» schwärmte Louis und wusste, mit Meo hatte er den Richtigen gefunden,
«und, na ja, ich habe nicht mehr allzu viel Geld übrig. Aber eine
Occasion-Taylor für einen Tausender wäre vielleicht zu bekommen. Ich bin
auch mit einer Einfacheren zufrieden, als Übergang.»

Louis zog die Schultern hoch. Meo wandte sich an Vanessa.

«Nessi, was denkst du zu «Gitarrentotal»? Da waren wir doch kürzlich
mal.»

«Ja, ich erinnere mich. Die haben viele Occasionen. Ist gleich um die
Ecke. Zürich-Wiedikon, oder?»

Sie schaute Louis an, als wolle sie prüfen, ob er sich hier auskenne.

«Wir müssen auch gleich weiter, Ciro, und können dich mitnehmen. Bist du
zu Fuss?» 

Meos Frage liess Louis die Augen aufreissen und verdutzt nicken.

Vollgepackt zirkelten die drei zwischen einzelnen Menschengruppen quer über den
Platz. Meo und Vanessa steuerten Louis drei Strassen weiter auf einen
weissen «VW-Bus California» hin.
Schon während des Gehens hatten sie lebhafte Gespräche. Die Chemie
stimmte.
Louis gelang, das Gespräch immer wieder von sich weg auf ihre
Geschichten zu lenken. Er fragte nach der Bedeutung von Meos Namen.

«Glaube, diese Frage verfolgt mich schon seit meiner Geburt,» Meo
seufzte gespielt, «meine Eltern sind beide Pädagogen und vor allem
meiner Mutter war es sehr wichtig, bei den eigenen Kindern nicht an ihre
Schulkinder erinnert zu werden. Also diese Namen seien schon vergeben,
sagte sie, besetzt. Und darum mussten sie bei der Namenssuche von
uns Kindern kreative Wege gehen,» er grinste, «aber Spass beiseite, mein Name
gefällt mir. Und manchmal erklär ich sogar gern, wie es dazu kam.»

Sie stiegen in den Bus, setzten sich vorne nebeneinander und Vanessa
startete den Motor. Louis registrierte überrascht das neuwertige,
zweifellos teure Fahrzeug. Die beiden Drehsitze neben dem Führersitz imponierten ihm. 
Vanessa gab das gesuchte Geschäft im Navi ein.

«Aemtlerstrasse 15, Wiedikon,» konzentriert studierte sie den
Stadtausschnitt, «hier, klar in der Nähe der Sihl,» dann
blickte sie auf die Uhr auf dem Armaturenbrett, «es reicht gut,
reinzuschauen, die schliessen erst in zwei Stunden.»

Vanessa fuhr los, um umgehend wieder anzuhalten. Die Rotlichter hier waren eng platziert.
Über den Weg wurde der Abendverkehr dichter und liess ihnen Zeit,
weiter zu plaudern. Louis stellte den beiden Frage um Frage. Sie
erzählten begeistert.
An der pädagogischen Hochschule kennengelernt, erklärte Vanessa, hätten
sie sich beide bereits nach vier Arbeitsjahren von der Schule
verabschiedet. Sie lebten nun von ihren Ersparnissen. Konnten sich umso
intensiver der Musik zuwenden und hatten mit privaten Musikstunden für
Kinder und Erwachsene etwas dazuverdient.
Meo schielte immer wieder zu Vanessa hinüber, pflichtete ihr
nachdrücklich bei oder ergänzte.
Louis hörte stumm zu. Die kamen ganz schön in Fahrt. Die Zeit würde
vergehen. Er musste nicht von sich sprechen. Und wenn alles gut ging,
kam er vielleicht zu einer neuen Gitarre.

«Privatschulen werden sich mehren, denk ich. Später vielleicht mal eine
Option für uns, nicht Nessi?» schloss Meo.

«Hm, eventuell. Jetzt sicher nicht.»

Mit einer abwehrenden Handbewegung steckte sie sich einen Kaugummi in
den Mund.

Der Bus sollte die beiden in ihr neues Leben führen. Im wahrsten Sinne, erzählte 
Meo. Sie würden ein Jahr lang die Welt erkunden. Wo immer
möglich als Strassenmusiker auftreten. Sie planten, ihr musikalisches
Können zu erweitern und ihre Geschäftsidee zu entwickeln. Eine
Kombination aus Auftritten, Strassenmusik und Instrumentalunterricht für
Kinder und Erwachsene sollte es werden. Und dies alles kostenminimal. 
Darum der Bus. Er sei als Reisegefährt, Wohnhaus, später als
Unterrichtsraum und Büro gedacht. Vanessa nickte immer wieder
beipflichtend, und Louis war, als würde sie nächstens abheben. 
Sie bebte vor Freude.

«Und das Beste, Ciro - » begann Vanessa anzufügen, als leise
\emph{«Scharlachrot»} erklang. Sie unterbrach sich, juckte hoch, und
drehte das Radio überlaut auf.

«Wow, hört, «Patent Ochsner», einfach ewig gut,» und schon sang
sie mit. Meo stimmte ein. Er nahm Vanessas Hand vom Steuer und hielt
sie. Sie fuhr einhändig weiter.
Trotz der quälenden Hitze hatte Louis Gänsehaut.
Beim Refrain überkam ihn eine tiefe Sehnsucht. Er kannte den Song nicht.
Den Berner-Dialekt des Sängers empfand er als angenehm und berührend. Er
blinzelte mehrmals, um seine Augen unter Kontrolle zu behalten. Meo zwinkerte zu Louis hinüber.

«Den hab ich schon als Kind gehört. Ein Lieblingssong meines Vaters.»

\qrblock{Patent Ochsner \emph{«Scharlachrot»}}{media/QR-Codes/021_Scharlachrot_Patent_Ochsner_Spotify.png}

Die letzten Takte wurden ausgeblendet. Ausgerechnet jetzt. Louis ärgerte
sich. Das Saxophon-Solo hatte es ihm angetan.
Die Stimme des Nachrichtensprechers löste Piep und Zeitangabe ab. Sechzehn Uhr.

«Und nun die Nachrichten des Tages.»

Vanessa drehte das Radio leiser.

«Sehr schön, nicht nur der Song, auch ihr zwei. Ihr singt verdammt gut.
Ich kannte die Band nicht. Schweizerdeutsche Songs schaffen es nicht bis
nach Italien.»

Dabei lachte Louis und die beiden stimmten mit ein.

«Ich stelle mal auf Italien um, Ciro, «Radio Rock», für dich. Den Sender
hören wir oft, nicht, Nessi?»

Vanessa nickte abwesend. 

«Oder findest du vielleicht «Il Sole 24 ORE», Meo? Ein abwechslungsreicher Sender.»

«Aber ja. Ich versuch`s.» 

Meo drehte am entsprechenden Knopf. 
Die Ampel stellte auf grün um und Vanessa fuhr an.
Auch hier plätscherten die Nachrichten dahin. Meo und Vanessa begannen
Louis «Patent Ochsner» mit blumigen Worten zu beschreiben. Sie redeten
abwechselnd. Bis Meo Luft holte und Louis hellhörig wurde.

«Moment mal, Meo, kannst du das Radio lauter stellen, italienische
Nachrichten?»

\sectionbreak

\indentedblock{
«…wie vor ungeklärt. Die sechsundzwanzigjährige Frau stürzte in
Mailand anlässlich einer Feier über die Felswand des «Roccia del Culto».
Die «Polizia di Stato Milano» ermittelt. Ob es sich um einen Unfall oder
ein Verbrechen handelt, wird weiterhin mit Hochdruck untersucht.
Gesucht werden Personen, die sich von Sonntag, dem 19. auf Montag, dem
20. August zwischen Mitternacht und 04:00 Uhr in direkter Umgebung des
«Roccia del Culto» aufgehalten und etwas beobachtet haben.
Über den Zustand des Opfers können aus ermittlungstechnischen Gründen
keine Angaben gemacht werden.
Sachdienliche Hinweise nimmt die «Polizia di Stato Milano» oder eine
andere Polizei-Dienststelle rund um die Uhr entgegen.»
Es folgte eine direkte Telefonnummer.}

\sectionbreak

Louis war, als höre er seinen Herzschlag von aussen. 
Jetzt war es soweit. Er wurde bestimmt gesucht. Warum nicht direkt mit
Namen, Beschreibung und allem Drum und Dran? 
Die verschwommene Sicht quälte ihn. Er nahm die Brille ab und drehte das
Gestell nachdenklich in den Händen herum.
Keine Angaben zu Giorgias Zustand, was das bedeutete? Er bekam kaum
Luft.
Heilfroh erstmal, dass Vanessa das Wort ergriff, sank er in sich
zusammen.

«Och, arme Frau,» sagte Vanessa, «Mailand, da kommst du doch her, Ciro?»

Sie blickte weiterhin geradeaus und Louis bemühte sich, neutral zu klingen.

«Ja, schlimm so was, hm…Kompliment, du verstehst italienisch,
Vanessa?»

Sie hob die Brauen, lächelte und Louis ergänzte, «danke, Meo, kannst wieder leiser stellen.»

«Verbrechen, oh, oh.»

Meo verzog den Mund und schaute belustigt zu Louis. Auch er schien die Meldung
verstanden zu haben.

«Ja, du, das wär mir was,» Louis lockerte sich, «gut, bin ich schon seit
Monaten unterwegs, sonst wer weiss.»

Er hob seine Hände und verdrehte die Augen.
Meo lachte schallend.

«Jetzt hab ich doch glatt die Höhe der Belohnung für Täterhinweise 
verpasst.»

Vanessa stoppte den Bus vor dem nächsten Rotlicht.

«Häh? Wurde gar keine ausgesetzt, und wofür denn?»

Louis spielte angestrengt mit.

«Musikalische Kopfgeldjäger, mal was Neues.»

Alle lachten. Das Thema beendet.

\sectionbreak

Was folgte, konnte Louis kaum fassen.
Es dauerte keine Stunde, bis er das Geschäft zusammen mit den beiden wieder verliess. 
Mit dieser sagenhaften Gitarre in einem
tiefschwarzen Koffer. Es war beinahe unheimlich. Ob er so viel Glück verdient hatte?
Draussen konnte er sich gerade noch halten. Beinahe wäre er an Meo hochgesprungen.

«Danke, danke, wow, das ist der Hammer, ich fass es nicht,» rief er begeistert.

Meo und Vanessa standen angetan daneben. Seine Unbändigkeit schien ihnen zu gefallen.
Louis konnte es nicht erwarten, sich hinzusetzen. Egal wo, auch hier,
direkt auf den heissen Vorplatz an der Hauptstrasse. Am liebsten hätte
er augenblicklich losgelegt, das Instrument ausprobiert, stundenlang, 
wie an Weihnachten, wenn er das ersehnte Spielzeug unter dem
Geschenkpapier hervorblitzen sah. Und dann ewig damit spielen durfte,
weil Maman an Weihnachten keine Zeit kannte.
Die zwei hatten ja keine Ahnung, zu was sie ihm verholfen hatten. Sie
retteten ihm gerade die letzte Facette seines Lebens.
Ganz selbstverständlich stieg Louis zu ihnen in den Bus.

«Wohin willst du, Ciro?»

Vanessa startete den Motor.
Louis blickte sie entgeistert an. Über die kommende Nacht hatte er noch
gar nicht nachgedacht. Seine Hände umklammerten den Hals des
Gitarrenkoffers zwischen seinen Beinen. Eine Hitzewelle überzog ihn.
Meo zog die Augenbrauen hoch.

«Erstmal Zeit für eine Session, ihr Lieben, kommst du zu uns? Und wo
übernachtest du eigentlich?»

Louis kramte nach einem Einfall.

«Werde was Günstiges suchen, eine Jugendherberge oder so.»

«Ach was, Ciro, du kannst bei uns übernachten. Wir sind heute auf dem
Vorplatz meiner Eltern stationiert. Ist gleich um die Ecke. Wir können
für dich das Vorzelt am Bus anbringen, ok?» 

Louis nickte erleichtert. Ja, es fühlte sich
definitiv an wie Weihnachten. Sie fuhren plaudernd nach Urdorf.

Kaum hatte Meo das Radio erneut aufgedreht, war Louis die Polizeimeldung
wieder in den Ohren. Keine Angaben zu Giorgias Zustand.

Sie waren angekommen.
